% Section 3.1, Wikipedia dataset. Two variants for supervisor choice (owner review 2026-07-08); appendix details at the bottom.
% ===== VARIANT 1 (main, conservative) =====
\textit{Wikipedia} is a corpus of 100 encyclopedic articles on ancient history
that we assemble for terminology-oriented evaluation. The articles are drawn
from the Russian edition across ten thematic sections (Sumer and Mesopotamia,
Ancient Egypt, Assyria, the Hittite Kingdom, Phoenicia, Achaemenid Iran,
Ancient India, Ancient China, Ancient Greece, and Ancient Rome), ten per
section. Within each section, candidates from the corresponding category
subtree are shuffled with a fixed seed and admitted through a deterministic
gate, which requires at least 30 unique in-text links and an earliest
associated Wikidata date before 500~CE, so the selection is reproducible and
involves no manual curation. We use this corpus primarily to evaluate
terminology extraction and Wikidata grounding. The hyperlinks that editors
place in article bodies act as human annotation: each link target resolves
deterministically to a Wikidata item, so every in-text link yields a gold
mention grounded to a QID at no additional labeling cost. This annotation is
precise but incomplete, since editorial convention is to link an entity only
at its first mention~\citep{damuel2023}, and our evaluation protocol accounts
for this (Section~\ref{sec:eval-metrics}). The same articles also serve as
source texts for the reference-free evaluation of translation and refinement
quality. After markup removal and paragraph segmentation, the corpus contains
2553 paragraphs and 160836 words, with an average paragraph length of 63 words
(Table~\ref{tab:data-statistics}).

% ===== VARIANT 2 (creative lead, recommended by orchestrator) =====
% \textit{Wikipedia} builds on an annotation resource that encyclopedia editors
% maintain by hand: in-text hyperlinks. Each link target resolves
% deterministically to a Wikidata item, so every link in an article body
% provides a gold entity mention grounded to a QID at no labeling cost, which
% makes the corpus a natural benchmark for terminology extraction and Wikidata
% grounding. We assemble 100 articles on ancient history from the Russian
% edition, spanning ten thematic sections (Sumer and Mesopotamia, Ancient Egypt,
% Assyria, the Hittite Kingdom, Phoenicia, Achaemenid Iran, Ancient India,
% Ancient China, Ancient Greece, and Ancient Rome). The selection is
% deterministic and involves no manual curation: within each section, candidates
% from the category subtree are shuffled with a fixed seed, and the first ten
% articles pass a gate requiring at least 30 unique in-text links and an
% earliest Wikidata date before 500~CE. The resulting annotation is precise but
% incomplete, as editors typically link an entity only at its first
% mention~\citep{damuel2023}, a property our evaluation protocol accounts for
% (Section~\ref{sec:eval-metrics}). Beyond grounding, the same articles serve as
% source texts for the reference-free evaluation of translation and refinement
% quality. The corpus comprises 2553 paragraphs and 160836 words, averaging 63
% words per paragraph (Table~\ref{tab:data-statistics}).

% ===== APPENDIX SNIPPET (construction details) =====
% \paragraph{Wikipedia corpus construction details.}
% Candidate pools are the article-namespace category subtrees of the ten section
% root categories. An article reachable from several sections is assigned to the
% first in a fixed order and used once. For the chronological criterion we take
% the earliest value among the Wikidata properties inception (P571), start time
% (P580), point in time (P585), date of birth (P569), and publication date
% (P577). Articles with no recorded date are retained. Articles typed as films,
% paintings, or museums are excluded. Shuffling uses seed 42. Article HTML was
% collected in July 2026. Because the Wikipedia category graph is noisy, roughly
% three of the 100 articles fall outside the target period. They are disclosed
% and retained to avoid manual curation.
